% =====================================================================
% 7. Conclusions (MLLB, Maria Martinez)
% =====================================================================
\section{Conclusions}\label{sec:mllb-conclusions}

\begin{enumerate}[leftmargin=1.4em,itemsep=2pt]
  \item \textbf{A shared pipeline makes architectures comparable.} One data pipeline, one training loop
        driven by a task object, one evaluation module and one log format served every classifier and the
        generators. Adding a model meant writing an \texttt{nn.Module}; nothing else changed, so the
        differences we report come from the models.
  \item \textbf{The silent errors are in the data and the shapes.} A wrong label rule, statistics computed
        on evaluation records, a wrong axis order or the wrong row of a bidirectional output all train
        without error and give plausible numbers. Each of them is now covered by a test that checks a known
        answer: published class counts, training-only statistics, element-by-element framing, the forward
        and backward final states, equality with \texttt{nn.MultiheadAttention}.
  \item \textbf{Writing attention by hand was worth it.} It made every step explainable (scores, scaling,
        mask, heads, $W_O$, Pre-LN), gave the attention maps for free, and the same module served
        Oriana's generator with a causal mask.
  \item \textbf{Reproducibility needs more than a seed.} Seeding the libraries, a seeded generator for the
        batch order, deterministic cuDNN and \texttt{CUBLAS\_WORKSPACE\_CONFIG} made the final runs of my
        blocks reproduce their validation runs exactly on the laptop's GPU.
  \item \textbf{Numbers from logs, comparisons on the same patients.} Generating every number from the
        logs, and comparing models on the same patient-bootstrap resamples seed by seed, is what allowed
        us to say which ablations mattered (reading all hidden states, the GRU, the convolutional front
        end, the heads and the residual path of self-attention) and which did not (the attention score, the
        number of heads beyond one block, the readout of the Transformer).
  \item \textbf{Limits.} Small models and budgets on a 4\,GB GPU, three seeds, and one dataset. The
        cost table shows the price of recurrence over long sequences: the LSTMs over 1000 raw steps, with and
        without attention, were the slowest configurations of my blocks (Table~\ref{tab:mllb-cost}).
\end{enumerate}

\subsection{My activities in the team}

Table~\ref{tab:contrib_mm_mllb} gives my role in each completed task, generated from the task log
(\texttt{docs/PLAN\_DL.md}); ``did (shared)'' marks joint tasks.

\begin{table}[htbp]
\centering
\caption{My role in each completed task (generated from \texttt{docs/PLAN\_DL.md}).}
\label{tab:contrib_mm_mllb}
% Generated by common/make_contributions.py from docs/PLAN_DL.md (MM). Do not edit.
{\small
\begin{tabularx}{\textwidth}{l>{\raggedright\arraybackslash}X>{\raggedright\arraybackslash}p{3.6cm}l}
\toprule
ID & Task & Role of Mar\'ia Mart\'inez & Status \\
\midrule
T0 & Dataset shortlist (8 candidates) and choice & did (shared) & done \\
T1 & \texttt{DECISIONS.md} and \texttt{PLAN\_DL.md} first version & did (shared), reviewed & done \\
T2 & Decide the transfer-learning approach: own pre-trained model (not confirmed with the professor; raise it in the class discussion) & did (shared) & done \\
T3 & Verify licence and recommended folds on PhysioNet & did (shared) & done \\
T4 & Local project structure, \texttt{.gitignore}, environment file; GitHub repository later & reviewed & done \\
T5 & Download script, label construction, fold split, \texttt{Dataset}/\texttt{DataLoader} & did & done \\
T6 & EDA: class balance, co-occurrence, demographics, signal quality, example traces & reviewed & done \\
T7 & Evaluation module: macro AUROC, bootstrap CI, JSON logging, \texttt{make\_numbers.py} & did & done \\
T8 & Shared training loop (zero\_grad, forward, loss, backward, step; train/eval modes) & did (shared) & done \\
T9 & Dense baseline + ablations & reviewed & done \\
T10 & CNN 1D + ablations & reviewed & done \\
T11 & LSTM/GRU (+bidirectional) and RNN + attention model + ablations & did & done \\
T12 & Self-attention step by step (single layer, attention maps) and Transformer + ablations & did & done \\
T13 & Transfer learning: ECG-GPT (T23) reused on the 5 superclasses, scratch vs frozen vs fine-tuning, 10\% and 100\% labels & did (shared) & done \\
T14 & Cross-cutting ablation: 12 leads vs single lead & did (shared) & done \\
T15 & Final comparison table and figures; comparison with published results & did (shared) & done \\
T16 & Literature: dataset and CNN-related papers & reviewed & done \\
T17 & Literature: RNN/Transformer-related papers & did & done \\
T18 & Reviewed first assignment (learning strategies) for the PRDL reports & did her own part, cross review & done \\
T19 & PRDL-DL report (theory + critical analysis), opening with a summary of at most 3 pages & did her own part & done \\
T20 & MLLB-DL report (Section 0 AI environment, code, training, evaluation), opening with a summary of at most 3 pages & did her own part & done \\
T21 & \texttt{AI\_USAGE\_LOG.md} and AI Error Bounty sheets & did her own part & done \\
T22 & Adversarial review before delivery; preparation of the class discussion and oral defence & planned: owner & todo \\
T23 & Generative autoregressive model (ECG-GPT): tokenisation of lead II, causal Transformer and LSTM, sampling with temperature, ablations (D13) & reviewed & done \\
T24 & Zero-shot learning: theory and a short experiment on the class CLIP notebook & did (shared) & done \\
T25 & Initial draft of the DL reports submitted (due 1 December 2026) & planned: owner & todo \\
\bottomrule
\end{tabularx}}

\end{table}

In code, my part was the data layer (\file{src/data.py}, \file{src/dataset.py}), the evaluation module
and the run logs (\file{src/metrics.py}, \file{src/compare.py}, \file{src/runlog.py}), the recurrent
models and the hand-written attention and Transformer (\file{src/models/rnn.py},
\file{src/models/attention.py}, \file{src/models/transformer.py}) with their experiment scripts and tests.
For each I went through the code and the tests with the assistant before launching the runs, read the
logs and checked the figures. The training loop and the area-2 experiments were joint work with Oriana,
whose code I reviewed; her implementations are summarised here from her block notes and credited to her.

